% Reconstructed text from V5 PDF page 7; non-native source
\par\smallskip\noindent\rule{\textwidth}{0.4pt}\par{\sffamily\footnotesize Original V5 page 7 -- recovered text}\par\begin{flushleft}\footnotesize\ttfamily
\par\smallskip
\par\smallskip
\noindent Diagnostics and Serving Infrastructure for MAS 7\par
\par\smallskip
\par\smallskip
\par\smallskip
\noindent systems surveyed, eviction mechanisms range from simple LRU (SGLang) to learned survival-\par
\noindent based predictions (the policy-driven runtime), with 8 systems incorporating agent-level meta-\par
\noindent data into their management decisions. The key distinction is between reactive eviction (trig-\par
\noindent gered by memory pressure) and predictive eviction (triggered by estimated future reuse), with\par
\noindent the latter providing more opportunity for diagnostic-aware retention.\par
\noindent  The distinctive question for this review is which state dependencies cross the diagnos-\par
\noindent tic/serving boundary, and which apparent similarities between the two layers do not constitute\par
\noindent resource competition. This question is formalized in Section 6.\par
\par\smallskip
\par\smallskip
\noindent 3 Review Methodology\par
\noindent We assembled a curated critical-review corpus through iterative database exploration, back-\par
\noindent ward/forward citation chaining, and eligibility-based selection. Because database-specific search\par
\noindent logs and screening decisions were not retained, we do not report a PRISMA-style numerical\par
\noindent flow or claim systematic completeness. The PRISMA 2020 guidelines [26] informed the re-\par
\noindent porting structure only. Figure 2 summarizes the sources, inclusion criteria, and final corpus\par
\noindent composition. We describe our search strategy, selection criteria, evidence-coding scheme, and\par
\noindent the classification taxonomy below.\par
\par\smallskip
\noindent 3.1 Search Strategy\par
\noindent We conducted iterative searches across five academic databases: IEEE Xplore, ACM Digital\par
\noindent Library, Scopus, Google Scholar, and arXiv. The initial search window spanned January 2024\par
\noindent through March 2026, followed by a targeted citation-chaining update through August 2026, and\par
\noindent a final update pass through November 2026 capturing systems announced at NeurIPS 2025,\par
\noindent ICML 2026, ACL 2026, EMNLP 2026, OSDI 2024, FAST 2025, and EuroSys 2025/2027. Be-\par
\noindent cause foundation-model multi-agent runtime research is rapidly evolving, we included preprints\par
\noindent alongside peer-reviewed publications when the described system contained reproducible imple-\par
\noindent mentation details. We prioritize work from 2024–2026 because foundation-model multi-agent\par
\noindent runtime research is unusually recent, while retaining older serving papers when they define\par
\noindent mechanisms directly reused by agentic systems.\par
\noindent  We constructed Boolean search queries combining terms from two domains:\par
\noindent  • Diagnostics: (“multi-agent” OR “LLM agent” OR “foundation model agent”) AND\par
\noindent  (“failure localization” OR “fault attribution” OR “debugging” OR “observability” OR\par
\noindent  “tracing” OR “guardrail” OR “safety intervention” OR “runtime monitoring”)\par
\noindent  • Infrastructure: (“multi-agent” OR “LLM serving” OR “agent serving”) AND (“KV\par
\noindent  cache” OR “memory management” OR “scheduling” OR “state sharing” OR “cache evic-\par
\noindent  tion” OR “prefix caching” OR “context reuse”)\par
\noindent  Cross-domain queries additionally combined diagnostic and infrastructure terms to capture\par
\noindent systems at the intersection (e.g., “KV cache eviction” AND “diagnostic evidence”). To make the\par
\noindent review more reproducible without overstating exhaustiveness, candidate papers were collected\par
\noindent iteratively from primary systems papers, adjacent surveys, major peer-reviewed venues, and\par
\noindent forward/backward citation tracing. Queries were adapted per database syntax. A paper enters\par
\noindent the core system tables only when its runtime role can be identified from the source and mapped\par
\noindent to one of the functions above. Because the corpus was not generated by a preregistered database\par
\noindent protocol with fixed search strings and PRISMA-style screening, the reported counts should be\par
\noindent interpreted as a curated landscape rather than an exhaustive census.\par
\end{flushleft}
